\documentclass[
  fontsize=10pt,
  twoside=false,
  secnumdepth=2,
  numbers=noenddot,
  fontmethod=modern
]{kaohandt}

\usepackage[american]{babel}
\usepackage{lecture-notes}
% Edit this file for each lecture.
\newcommand{\CourseHeader}{LEARN ML BY BUILDING}
\newcommand{\LectureNumber}{NN}
\newcommand{\LectureTitle}{Lecture Notes Template}
\newcommand{\LectureSubtitle}{Problem-first structure for concise technical lectures}
\newcommand{\LecturePdfSubject}{Problem formulation, assumptions, method, and evaluation}
\newcommand{\LectureAbstract}{A reusable visual and editorial system for turning current handwritten notes, supporting transcript context, and verified course artifacts into focused technical lecture notes.}
\newcommand{\LectureAcknowledgmentText}{These notes are based on the instructor's handwritten notes and transcripts of lecture discussions and were editorially polished and typeset with assistance from OpenAI Codex and Anthropic Claude. The instructor reviewed and is responsible for the final content.}


\titlehead{\textcolor{LectureTeal}{\sffamily\bfseries\small \CourseHeader}}
\subject{LECTURE \LectureNumber}
\title[\LectureTitle]{\textcolor{LectureInk}{\LectureTitle}}
\subtitle{\LectureSubtitle}
\author[Learn ML by Building]{}
\date{}

\hypersetup{pdfsubject={\LecturePdfSubject}}

\begin{document}

\maketitle
\pagestyle{pagenum.scrheadings}
\thispagestyle{pagenum.scrheadings}

\begin{abstract}
\noindent
\LectureAbstract
\end{abstract}

\begin{widepar}
\begin{learningbox}
\textbf{Learning goals.} By the end of this lecture, you should be able to
\begin{enumerate}[label=\textcolor{LectureTeal}{\arabic*.}, leftmargin=1.8em, itemsep=0.15em]
  % Replace these with goals supported by the material actually taught.
  \item formulate the practical problem and the evidence needed to evaluate it;
  \item explain the method and the assumption that makes it plausible; and
  \item diagnose a failure at the correct layer of the system.
\end{enumerate}
\end{learningbox}
\end{widepar}

\section{Begin with the problem}

Open with the decision or failure that motivates the lecture. Introduce the algorithm only after the problem has a clear formulation, and keep the opening example proportional to its role in the lecture.

\marginidea{Use the margin for a brief design implication, not a second narrative thread.}

\subsection{Move from the problem to a formulation}

Show the conceptual translation explicitly:

\[
\begin{gathered}
\text{practical problem}\;\longrightarrow\;\text{mathematical condition}\\[-1mm]
\longrightarrow\;\text{model or objective}\;\longrightarrow\;\text{evaluation evidence}.
\end{gathered}
\]

Define notation immediately before it is used and explain why each component matters. One focused worked example is usually more valuable than several near-duplicates.

\begin{widepar}
\centering
\begin{tikzpicture}[node distance=4mm, process/.append style={font=\sffamily\footnotesize, inner xsep=2mm}]
  \node[process, fill=LectureTealLight, draw=LectureTeal, text width=20mm] (problem) {\textbf{Problem}\\real decision};
  \node[process, right=of problem, text width=22mm] (formulation) {\textbf{Formulation}\\data and target};
  \node[process, right=of formulation, fill=LectureBlueLight, draw=LectureBlue, text width=21mm] (assumption) {\textbf{Assumption}\\why it may work};
  \node[process, right=of assumption, fill=LectureGoldLight, draw=LectureGold, text width=19mm] (method) {\textbf{Method}\\learn or predict};
  \node[process, right=of method, fill=LecturePurpleLight, draw=LecturePurple, text width=20mm] (evidence) {\textbf{Evidence}\\evaluate outcomes};

  \draw[strong link] (problem)--(formulation);
  \draw[strong link] (formulation)--(assumption);
  \draw[strong link] (assumption)--(method);
  \draw[strong link] (method)--(evidence);

  \draw[-{Stealth[length=2.5mm]}, draw=LectureCoral, line width=1.2pt, rounded corners=2mm]
    (evidence.south) -- ++(0,-10mm) -| node[pos=0.52, below, font=\sffamily\footnotesize, text=LectureCoral]{diagnose and revise the right layer} (problem.south);
\end{tikzpicture}

\FigureCaption{fig:template-pipeline}{The standard figure format: native TikZ, shared colors and node styles, a compact title-free composition, and a caption that states the message.}
\end{widepar}

\subsection{State assumptions and diagnose them}

Do not treat a model assumption as an isolated mathematical sentence. Trace it through the system that creates the model's inputs and evaluates its outputs.

\begin{cautionbox}[Use emphasis selectively]
Reserve colored boxes for a limitation, warning, or worked result that students should genuinely notice. Keep ordinary principles in ordinary prose.
\end{cautionbox}

When several layers must be compared exactly, use a LaTeX table:

\begin{widepar}
\footnotesize
\begin{tabularx}{\linewidth}{L{0.18\linewidth} L{0.29\linewidth} Y Y}
\toprule
\tablehead{Layer} & \tablehead{Assumption} & \tablehead{Evidence} & \tablehead{Possible response} \\
\midrule
Problem & The target represents the real decision & Stakeholder definition and downstream outcome & Reframe the target \\
Representation & Available features preserve relevant information & Data audit and counterexamples & Revise collection or features \\
Model & The method's structural assumption is useful & Held-out behavior and error analysis & Change the ruler or method \\
Evaluation & The test reproduces future use & Prospective or deployment-aligned evidence & Redesign the validation \\
\bottomrule
\end{tabularx}
\TableCaption{tab:template-diagnosis}{The standard table format: native LaTeX, restrained rules, semantic columns, and a concise interpretive caption.}
\end{widepar}

\begin{widepar}
\medskip
{\small
\noindent\textcolor{LectureTeal}{\sffamily\bfseries Lecture summary.}
Replace this with one sentence that connects the practical problem, the model's main assumption, and the evidence used to diagnose success or failure.

\smallskip
\noindent\textcolor{LectureTeal}{\sffamily\bfseries Further reading.}
\begin{itemize}[leftmargin=1.35em, itemsep=0.05em, topsep=0.2em]\raggedright
  \item the primary paper or canonical reference;
  \item official technical documentation; and
  \item the corresponding course notebook or project.
\end{itemize}
}

\lectureacknowledgment{\LectureAcknowledgmentText}
\end{widepar}


\end{document}
